\scanpage{135}
% archive sheet number: 134

en fait, comme on a

\[
H\unsure{\longrightarrow|-}L,\quad H'\unsure{\longrightarrow|-} L'
\]

on en déduit

\[
P\overset{\mathrm{dfn}}{=}H\otimes H'\longrightarrow L\otimes L'\qquad\text{épi}
\]

nul sur $1\otimes H'$ et sur $H\otimes1$, donc sur
$B$, et on a la suite exacte

\diagram{\begin{tikzcd}[ampersand replacement=\&,column sep=small]
0\arrow[r] \& B\arrow[r] \& P\arrow[r] \& \deleted{\unsure{1}}L\otimes L'\arrow[r] \& 0
\end{tikzcd}}

\deleted{Ceci dit}
Ainsi les formes linéaires \inserted{$T_P$} sur
$P$ dont la restr. à $B$ est $T_B$, font
un torseur \inserted{$\underline L_P$} sous $(L\otimes L')^\vee$, L'application
can.

% slanted
\margin{i.e. les formes bil. $H\times H'\xrightarrow{f}k$, telles que $f(u,1)=\mathrm{Tr}_H(u)$, $f(1,v)=\mathrm{Tr}_{H'}(v)$}

\[
\begin{aligned}
\zeta:\ \underline{\Phi}_B&\longrightarrow\underline L_P\\
\varphi_B&\longmapsto (x\otimes y\longmapsto\varphi_B(x\deleted{\unsure{\otimes1,1\otimes{}}},y))
\end{aligned}
\]

est compatible avec les opérations
de $(L\otimes L')^\vee$. Donc $\zeta$ est iso.

Notons le

\uline{Lemme} Soient $H,H'$ deux modules vectoriels
de rang 2 sur $k$, $P\simeq H\otimes H'$. On définit
une fonction quadratique canonique

\diagram{\begin{tikzcd}[ampersand replacement=\&]
P\arrow[r,"{\unsure{\delta_{H,H'}|q_{H,H'}|\Delta_{H,H'}}}"] \& \underbrace{\det H}_{L}\otimes\underbrace{\det H'}_{L'}
\end{tikzcd}}

en interprétant $P\simeq\mathrm{Hom}(H^\vee,H')\longrightarrow\mathrm{Hom}(\det H^\vee,\det\unsure{H'})$

